\section{Cloudfront - Part 1}
\begin{itemize}
    \item Content Delivery Network (CDN)
    \item Improves read performance, content is cached at the edge
    \item 225+ Point of Presence globally (215+ Edge Locations and 13 Regional Edge Caches)
    \item Protect against Network and Application layer attacks (e.g DDoS attacks)
    \item Integration with AWS Shield, AWS WAF and route 53
    \item Can expose external HTTPS and can talk to internal HTTPS backends
    \item Supports WebSocket protocol
\end{itemize}

\subsection{CloudFront - Origins}
\begin{itemize}
    \item S3 Bucket
    \item For distributing files
    \item For uploading to S3 (using CloudFront as an ingres)
    \item Enhanced security with CloudFront Origin Access Control (OAC)
    \item MediaStore Container and MediaPackage Endpoint
    \item To deliverVideo on Demand (VOD) or live streaming video using AWS Media Services
    \item VPC Origin
    \item For applications hosted in VPC private subnets
    \item Application Load Balance / Network Load Balancer / EC2 Instances
    \item Customer Origin (HTTP)
    \item API Gateway (for more control.... otherwise use API Gateway Edge)
    \item S3 Bucket configured as a website (enable Static Website hosting)
    \item Any HTTP backend you want
\end{itemize}

\subsection{CloudFront - S3 as an Origin}
%Diagram

\subsection{CloudFront vs S3 Cross Region Replication}
\begin{itemize}
    \item CloudFront:
    \begin{itemize}
        \item Global Edge network
        \item Files are cached for a TTL (maybe a day)
        \item Great for static content that must be available everywhere
    \end{itemize}
    \item S3 Cross Region Replication
    \begin{itemize}
        \item \item Must be setup for each region you want replication to happen
        \item Files are updated in near real-time
        \item Read only
        \item Great for dynamic content that needs to be available at low-latency in few regions
    \end{itemize}
\end{itemize}

\subsection{CloudFront - ALB or EC2 as an origin Using VPC Origins}
\begin{itemize}
    \item Allows you to deliver content from you applications hosted in your VPC private subnets (no need to expose them on the Internet)
    \item Deliver traffic to private
    \item Application Load Balancer
    \item Network Load Balancer
    \item EC2 Instances
\end{itemize}

%Diagram

\subsection{CloudFront - EC2 or ALB as an origin}
%Diagram

\subsection{CloudFront - Restrict Access to Application Load Balancers and Custom Origins}
\begin{itemize}
    \item Prevent direct access to you ALB or Custom Origins (only access through CloudFront)
    \item First, configurations CloudFront to add a CustomHTTPHeader to requests it sends to the ALB
    \item Second, configure the ALB to only forward requests that contain that Customer HTTP Header
    \item Keep the custom header name and value secret!
\end{itemize}
%Diagram

\subsection{CloudFront - Origin Groups}
\begin{itemize}
    \item To increase high-availability and do failover
    \item Origin Group: one primary and on secondary
    \item If the primary origin fail; the second on is used
    \item Origins can be cross AWS regions
\end{itemize}

%diagram


\section{Cloudfront - Part 2}

\subsection{CloudFront Geo Restrictions}
\begin{itemize}
    \item - You can restrict who can access your distribution
    \item Allow list: Allow you users to access your content only if they're in one of the countries on a list of approved
    \item Block list: Prevent your users from accessing your content if they're in one of the countries on a blacklist of banned countries..... lol
    \item The "country" is determined using a third party Geo-IP database
    \item Use case: Copyright Laws to control access to content
    \item Note: the geo header CloudFront-Viewer-Country is in Lambda at Edge
\end{itemize}

\subsection{CloudFront - Pricing}
\begin{itemize}
    \item CloudFront Edge locations are all around the world
    \item The cost of data out per edge location varies
\end{itemize}
%Insert table

\subsection{CloudFront - Price Classes}
\begin{itemize}
    \item You can reduce the number of edge locations for cost reductions
    \item Three price classes
    \item Price Class All: all regions - best performance
    \item Price Class 200: most regions, but excludes the most expensive regions
    \item Price Class 100: only the least expensive regions
\end{itemize}
%digrams

%prety map diagram

\subsection{CloudFront Signed URL Diagram}
\begin{itemize}
    \item Signed URL with expiration to control access to content in CloudFront
    \item The Signed URL are generated by an API call into CloudFront as a trusted signer
\end{itemize}
%Diagram

\subsection{CloudFront Signed URL vs S3 Pre-Signed URL}
\begin{itemize}
    \item CloudFront Signed URL:
    \item Allow access to a path, no matter the origin
    \item Account wide key-pair, only the root can manage
    \item Can filter by IP, path, date, expiration
    \item Can leverage caching features
\end{itemize}

\begin{itemize}
    \item - S3 Pre-Signed URL:
    \item Issue a request as the person who pre-signed URL
    \item Uses the IAM key of the signed IAM principal
    \item Limited lifetime
\end{itemize}

%Diagram

\subsection{CloudFront - Custom Error Pages}
\item Return an object to the viewer (e.g html) when your origin returns an HTTP 4xx or 5xx status code to CloudFront
\item Use Error Caching Minimum TTL to specify how long CloudFront caches the customer error pages

%diagram


\section{Lambda at Edge and CloudFront Functions}

\subsection{CloudFront - Customisation at the edge}
\begin{itemize}
    \item Many modern application execute some form of the logic at the edge
    \item Edge function:
    \item A code that you write an attach to CloudFront distributions
    \item Runs close to your users to minimise latency
    \item Doesn't have any cache, only to change requests/ responses
    \item CloudFront provides two types: CloudFront Functions and Lambda and Edge
\end{itemize}
Use cases:
\begin{itemize}
    \item Manipulate HTTP requests and responses
    \item Implement request filtering before reaching your application
    \item User authentication and authorisation
    \item Generate HTTP resources at the edge
    \item A/B Testing
    \item Bot mitigation
    \item You don't have to manage any servers, deployed globally
\end{itemize}

\subsection{CloudFront Functions and Lamabda at Edge}
%digram

\subsection{CloudFront - CloudFront Functions}
\begin{itemize}
    \item Lightweight functions written in JavaScript
    \item For high\itemscale latency-sensitive CDN customisation
    \item Sub-ms startup times, million of requests per second
    \item RuN at Edge Locations
    \item Processed-Based isolation
    \item Used to changeViewer requests and responses:
    \item Viewer Requests: after CloudFront recevives a reqeusts from a viewer
    \item Viewer Response: before CloudFront forwards the response to the viewer
    \item Native feature of CloudFront (manage code entirely within CloudFront)
\end{itemize}
%diagram

\subsection{CloudFront - Lambda at Edge}
\begin{itemize}
    \item - Lambda functions written in NodeJS or Python
    \item Scala to 1000s of reqeust/ second
    \item Runs at the nearest Regional Edge Cache
    \item VM-based isolation
    \item Used to change CloudFront requests and responses
    \item Viewer Request - after CloudFront recieves a reqeust from a viewer
    \item Origin Request - before CloudFront forwards the response to the origin
    \item Orign response - after CloudFront recieves the response from the origin
    \item Viewer Response - before CloudFront forwards the response to the viewer
    \item Author your functions in one AWS Region (some region) then CloudFront replicats to its locations
\end{itemize}

\subsection{CloudFront Functions with Lambda at Edge}
CloudFront Functions and Lambda at Edge can be used together
%digram
Note: You can't combine CloudFront Functions and Lambda at Edge in viewer events ( viewer reqeust and viewer response)

\subsection{Using Lambda at Edge only}
Use when you need some of the capabilities of Lambda at Edge that aren't available with CloudFront Functions (e.g longer execution time, network access,....)
%Diagram here

\subsection{CloudFront Funtions vs. Lamda at Edge}
%table here

\subsection{CloudFront Functions vs Lambda@Edge Use Cases}

\begin{itemize}
    \item CloudFront Functions
    \begin{itemize}
        \item Cache key normalization
        \item Transform request attributes (headers, cookies, query strings, URLs) to create an optimal cache key
        \item Header manipulation
        \item Insert, modify, or delete HTTP headers in requests or responses
        \item URL rewrites and redirects
        \item Request authentication and authorization
        \item Create and validate user-generated tokens (e.g., JWTs) to allow or deny requests
    \end{itemize}

    \item Lambda@Edge
    \begin{itemize}
        \item Longer execution time (up to several seconds)
        \item Adjustable CPU and memory allocation
        \item Support for third-party libraries and dependencies
        \item Access to AWS services through the AWS SDK
        \item Network access to external services for processing
        \item Access to the request body
        \item File system access via the Lambda execution environment
        \item Suitable for complex request and response processing
    \end{itemize}
\end{itemize}

\subsection{CloudFront Functions vs. Lambda at Edge - Authentication and Authorisation}
CloudFront Functions
%diagram

Lambda at Edge
%diagram

\subsection{Lambda@Edge: Loading content based on User-Agent}
%diagram here

\subsection{Lambda at Edge - Global Application}
%diagram


\section{Lambda at Edge Reduce Latency}

\subsection{Lambda at Edge - Route to different origin}
%diagram


\section{Amazon ElastiCache}

\subsection{Amazon ElasticCache Overview}
\begin{itemize}
    \item The same way RDS is to get managed Relational Databases.....
    \item ElastiCache is to get managed Redis or Memcached
    \item Caches are in memory databases with really high performance and low latency
    \item Helps reduce load off of databases for read intensive workloads
    \item Helps make you application stateless
    \item AWS takes care of OS maintenance / patching, optimisations, setup, configuration, monitoring failure recovery and backups.
    \item \textbf{Using ElastiCache involves heavy application code changes}
\end{itemize}

\subsection{ElastiCache Solution Architecture - DB Cache}
\begin{itemize}
    \item Applications queries ElastiCache, if not available, get from RDS and store in ElasticCache
    \item Helps relive load in RDS
    \item Cache must have an invalidation strategy to make sure only the most current data is used there.
\end{itemize}
%diagram

\subsection{ElastiCache Solution Architecture - User Session Store}
\begin{itemize}
    \item User logs into any of the application
    \item The application writes the session data into ElastiCache
    \item The user hits another instance of our application
    \item The instance retrieves the data the user is already logged in
\end{itemize}
%diagram

\subsection{ElastiCache - Redis vs Memcached}
\begin{itemize}
    \item - Redis
    \item Multi AZ with Auto-Failover
    \item Read Replicas to scale reads and have high availability
    \item Persistent, Data Durability: Append Only File (AOF), backup and restore features
    \item Memcached
    \item Multi-node for partitioning of data (sharding)
    \item Non Persistent
    \item Backup and restore (Serverless)
    \item Multi-threaded architecture
\end{itemize}


\section{Handling Extreme Rates}
%TODO diagram
